\section{Innledning}
    Denne oppgaven tar for seg DC- og AC-analyse og forståelse av felles emitter-forsterkere, som er en grunnleggende kobling innen analog elektronikk. Formålet er å undersøke hvordan en slik forsterker påvirker signaler, med særlig fokus på forsterkning, faseforsyvning samt ingangs- og utgangsimpedans.
    Gjennom oppgaven vil det bli analysert hvordan små variasjoner i inngangssignalet gir endringer i utgangssignalet, og hvordan kretsens oppbygning påvirker dette forholdet.

    Oppgaven er avgrenset til kun det tilgjengelige måleutstyret og komponentene som er tilgjengelig på laboratoriet.

    Hensikten med oppgaven er å øke forståelsen om hvordan felles emitter-forsterker fungerer, samt kunne forklare sentrale begreper knyttet til denne typen forsterkerkretser. 
